\chapter{Algèbres quadratiques d'ordre \(9\)}
\label{chap:algebras-cuadraticas-orden-nueve}
\index{algèbre quadratique}
\index{corps fini \(\mathbb F_9\)}

L'identité \(A_W^2=-I_6\), obtenue dans la section précédente par réduction
de la matrice de conférence de Paley, contient davantage d'information qu'une
orthogonalité modulaire. Elle détermine une extension des scalaires de degré deux,
transforme l'espace ternaire de dimension six en un espace de dimension
trois sur le corps à neuf éléments et permet d'écrire le code de
Golay ternaire comme le graphe d'une unité quadratique. Cette section développe
ces conséquences et les distingue de l'arithmétique de caractéristique neuf
utilisée par la branche de Hensel.

\section{Structure de \texorpdfstring{\(\mathbb F_9\)}{F9} induite par Paley}

Le polynôme \(X^2+1\) n'a pas de racines dans \(\mathbb F_3\) : ses valeurs en
\(0,1,2\) sont \(1,2,2\). Il est donc irréductible, et le quotient
\[
 \mathbb F_9
 :=\mathbb F_3[\iota]/(\iota^2+1)
\]
est un corps à neuf éléments.

\begin{theorem}[Structure quadratique induite par \(A_W\)]
\label{thm:f9-tres-desde-paley}
Soit \(V=\mathbb F_3^6\), écrit au moyen de vecteurs lignes, et soit
\(J(v)=vA_W\). La règle
\[
 (a+b\iota)v:=av+bJ(v),
 \qquad a,b\in\mathbb F_3,
\]
fait de \(V\) un espace vectoriel de dimension trois sur
\(\mathbb F_9\). Dans les coordonnées fixées par la construction de Paley,
\[
 \mathbb F_3^6\cong\mathbb F_9^3.
\]
\end{theorem}

\begin{proof}
L'égalité \(J^2=-I\) fait que l'évaluation
\(\mathbb F_3[X]\to\operatorname{End}_{\mathbb F_3}(V)\), \(X\mapsto J\),
s'annule sur l'idéal \((X^2+1)\). L'action se factorise donc par
le corps \(\mathbb F_9\). Comme
\([\,\mathbb F_9:\mathbb F_3\,]=2\) et
\(\dim_{\mathbb F_3}V=6\), on obtient
\(\dim_{\mathbb F_9}V=3\).
\end{proof}

L'isomorphisme dépend du système de coordonnées qui fixe \(A_W\). L'énoncé
ne se déduit pas du seul cardinal \(3^6=9^3\) : la matrice exhibe la
multiplication par tous les scalaires \(a+b\iota\) et fournit donc
la structure linéaire complète.

\section{Le code ternaire comme graphe isotrope}

L'application génératrice du code est
\[
 c(q)=(q,qA_W).
\]
Après identification de \(J\) à la multiplication par \(\iota\), son image
admet l'écriture
\[
 C_W=\{(q,\iota q):q\in\mathbb F_9^3\},
\]
où les deux composantes sont considérées, par restriction des scalaires, comme des
vecteurs de \(\mathbb F_3^6\).

\begin{proposition}[Maximalité isotrope]
\label{prop:golay-grafo-isotropico}
Par rapport à la forme bilinéaire standard de \(\mathbb F_3^{12}\), le
sous-espace \(C_W\) est totalement isotrope et maximal parmi les sous-espaces
totalement isotropes.
\end{proposition}

\begin{proof}
La matrice génératrice \(G=(I_6\mid A_W)\) satisfait
\[
 GG^{\mathsf T}=I_6+A_WA_W^{\mathsf T}=I_6+2I_6=0.
\]
L'espace engendré est donc totalement isotrope. Il est de
dimension six, la moitié de la dimension de l'espace ambiant non dégénéré,
et est donc maximal ; en fait, il coïncide avec son dual orthogonal.
\end{proof}

La racine de \(-I\) apparaît ainsi comme une transformation finie d'incidence.
On n'introduit pas le corps complexe analytique dans le système : on obtient une
structure quadratique sur \(\mathbb F_3\), suffisante pour organiser la
polarité, le code et ses deux espaces propres après extension des
scalaires.

\section{Caractéristique \(9\) et caractéristique \(3\)}

Deux ensembles de \(729\) états interviennent dans la construction :
\[
 (\mathbb Z/9\mathbb Z)^3
 \qquad\text{et}\qquad
 \mathbb F_9^3.
\]
Ils ne sont isomorphes ni comme groupes additifs ni comme anneaux. Le premier est de
caractéristique neuf et contient des éléments d'ordre neuf ; le second est de
caractéristique trois et tout son groupe additif est d'exposant trois.

Pour \(a,b\in\mathbb F_3\), en choisissant des représentants
\(\widetilde a,\widetilde b\in\{0,1,2\}\), soit
\[
 \beta(a,b)=\widetilde a+3\widetilde b\pmod 9.
\]
L'application en coordonnées induite par \(\beta\) est une bijection
ensembliste entre \(\mathbb F_3^6\) et \((\mathbb Z/9\mathbb Z)^3\), mais ne
respecte pas l'addition. En effet, si
\[
 \kappa(a,c)
 :=\left\lfloor
 \frac{\widetilde a+\widetilde c}{3}
 \right\rfloor,
\]
alors
\[
 \beta(a,b)+\beta(c,d)
 =\beta\bigl(a+c,\ b+d+\kappa(a,c)\bigr),
\]
avec les coordonnées du second membre réduites modulo trois. Par exemple,
\[
 \beta(2,0)+\beta(2,0)=\beta(1,1),
\]
tandis que la somme en coordonnées dans \(\mathbb F_3^2\) donnerait \((1,0)\). La
seconde coordonnée de \(\beta(1,1)\) conserve le quotient de la division
euclidienne ; ce terme est précisément la donnée qui disparaît lors de la
réduction composante par composante.

Cette distinction type les deux lectures du cardinal \(729\). La réalisation
\((\mathbb Z/9\mathbb Z)^3\) conserve le reste et le quotient ; la réalisation
\(\mathbb F_9^3\) transforme chaque paire ternaire en un scalaire d'une extension
quadratique. L'une prolonge l'arithmétique avec mémoire de division ; l'autre
organise l'incidence exceptionnelle.

\begin{proposition}[Deux lectures typées du même support ternaire]
\label{prop:dos-lecturas-tipadas-729}
Une fois fixées la section ensembliste utilisée dans \(\beta\), la base, la
polarisation et l'opérateur \(J\) du
\cref{thm:f9-tres-desde-paley}, il existe le diagramme de bijections ensemblistes
\[
\boxed{
(\mathbb Z/9\mathbb Z)^3
\xleftarrow{\;\beta^{\times3}\;}
\mathbb F_3^6
\xrightarrow{\;\iota_J\;}
\mathbb F_9^3.}
\]
La flèche de gauche transporte la somme nonadique tordue par le cocycle de
retenue \(\kappa\) ; la flèche de droite transporte la structure linéaire sur
\(\mathbb F_9\) déterminée par \(J^2=-I\). Les deux flèches partagent le
support en coordonnées \(\mathbb F_3^6\), mais n'identifient pas leurs lois
algébriques.
\end{proposition}

\begin{proof}
La formule
\[
 \beta(a,b)+\beta(c,d)
 =\beta\bigl(a+c,b+d+\kappa(a,c)\bigr)
\]
prouve que \(\beta^{\times3}\) est une bijection avec somme tordue et n'est pas
un isomorphisme à partir de la somme composante par composante de
\(\mathbb F_3^6\). Par le \cref{thm:f9-tres-desde-paley}, le choix de
\(J\) transforme ce même espace ternaire de dimension six en un espace
de dimension trois sur
\(\mathbb F_3[J]\cong\mathbb F_9\), ce qui définit \(\iota_J\).
\end{proof}

La proposition est une synthèse démontrée avec une structure de départ
explicite. Elle n'affirme pas une identification canonique entre
\((\mathbb Z/9\mathbb Z)^3\) et \(\mathbb F_9^3\) : le premier contient des
éléments additifs d'ordre neuf, tandis que le second possède un exposant
additif trois. Le cocycle mesure précisément l'information qui serait perdue en
confondant les deux lois.

\section{Classification des algèbres quadratiques}

Sur un corps de caractéristique différente de deux, une algèbre quadratique
unitaire se classe selon le type de son polynôme de définition. Sur
\(\mathbb F_3\), on peut choisir les trois représentants
\[
 \mathcal Q_{\mathrm e}
 =\mathbb F_3[t]/(t^2+1)\cong\mathbb F_9,
\]
\[
 \mathcal Q_{\mathrm n}
 =\mathbb F_3[\varepsilon]/(\varepsilon^2),
 \qquad
 \mathcal Q_{\mathrm s}
 =\mathbb F_3[r]/(r^2-1)
 \cong\mathbb F_3\times\mathbb F_3.
\]

\begin{theorem}[Classification des algèbres quadratiques sur \(\mathbb F_3\)]
\label{thm:clasificacion-cuadratica-nonadica}
Les relations
\[
 J^2=-I,\qquad N^2=0,\qquad R^2=I
\]
définissent des actions de
\(\mathcal Q_{\mathrm e}\), \(\mathcal Q_{\mathrm n}\) et
\(\mathcal Q_{\mathrm s}\), respectivement. Ces trois algèbres possèdent
neuf éléments et sont deux à deux non isomorphes : la première est un corps ; la
deuxième contient des nilpotents non nuls ; la troisième contient des idempotents non
triviaux.
\end{theorem}

\begin{proof}
Chaque relation fait se factoriser l'évaluation de
\(\mathbb F_3[t]\) par l'idéal correspondant. L'irréductibilité de
\(t^2+1\) prouve que \(\mathcal Q_{\mathrm e}\) est un corps. Dans
\(\mathcal Q_{\mathrm n}\), la classe \(\varepsilon\) est non nulle et satisfait
\(\varepsilon^2=0\). Dans \(\mathcal Q_{\mathrm s}\), comme \(2^{-1}=2\) dans
\(\mathbb F_3\), les éléments
\[
 e_+=2(1+r),\qquad e_-=2(1-r)
\]
sont des idempotents orthogonaux non triviaux dont la somme est l'unité. Être un corps,
posséder un nilpotent non nul et posséder des idempotents non triviaux sont des
propriétés invariantes par isomorphisme.
\end{proof}

Dans trois copies de la représentation régulière de dimension deux, on obtient
trois actions de dimension six sur \(\mathbb F_3\). La relation
\(J^2=-I\) correspond à la polarité de Paley ; la relation \(N^2=0\)
décrit la linéarisation d'une extension avec quotient ; la relation
\(R^2=I\) sépare deux termes au moyen des idempotents \(e_\pm\). Cette
classification fournit le cadre algébrique fini dans lequel les
structures elliptique, non réduite et scindée se distinguent sans faire appel à une
géométrie continue.

La construction ne remplace pas \(\mathbb C\) : le corps \(\mathbb F_9\) ne
possède ni la topologie, ni la complétude, ni la structure analytique complexe. Le
résultat exact est autre : une racine carrée de \(-I\) et l'extension
des scalaires correspondante émergent de la matrice d'incidence déjà construite
au sein du système ternaire.

\section{Source quadratique entière et spécialisations finie et réelle}
\label{sec:fuente-cuadratica-integral-rev3}

La racine finie et la racine opératorielle réelle se coordonnent au moyen d'une
présentation entière commune. Soit
\begin{equation}
 \mathcal Q_{\mathbb Z}
 :=\mathbb Z[u]/(u^2+1).
 \label{eq:fuente-cuadratica-integral-rev3}
\end{equation}
Ses changements de base sont
\[
 \mathcal Q_{\mathbb Z}\otimes_{\mathbb Z}\mathbb F_3
 \cong\mathbb F_3[u]/(u^2+1)
 \cong\mathbb F_9,
\]
et
\[
 \mathcal Q_{\mathbb Z}\otimes_{\mathbb Z}\mathbb R
 \cong\mathbb R[u]/(u^2+1).
\]

\begin{theorem}[Deux réalisations d'une présentation quadratique entière]
\label{thm:dos-realizaciones-fuente-cuadratica-rev3}
La représentation finie \(u\mapsto A_W\) sur
\(\mathbb F_3^6\) et la représentation réelle \(u\mapsto J_{\mathrm e}\) sur
\(V_{\mathrm e}\), construite dans le
\cref{thm:tres-generadores-concretos-trit}, réalisent, après leurs
changements de base respectifs, la même relation entière
\[
 u^2+1=0.
\]
\end{theorem}

\begin{proof}
L'égalité \(A_W^2+I_6=0\) fait se factoriser l'évaluation
\(\mathbb F_3[u]\to\operatorname{End}_{\mathbb F_3}(\mathbb F_3^6)\) par
\((u^2+1)\). De même, \(J_{\mathrm e}^2+I=0\) fait se factoriser
\(\mathbb R[u]\to\operatorname{End}_{\mathbb R}(V_{\mathrm e})\). Les deux
représentations proviennent de spécialisations de
\(\mathcal Q_{\mathbb Z}\) ; il n'existe et n'est utilisé aucun homomorphisme unitaire
\(\mathbb F_9\to\mathbb R\).
\end{proof}

La source commune coordonne des relations, elle n'identifie pas des modules de
caractéristiques distinctes. La branche finie conserve l'incidence ; la branche réelle
incorpore la structure analytique nécessaire à l'exponentielle opératorielle.

\section{\(3\) générateurs concrets des régimes TRIT}
\label{sec:tres-generadores-concretos-trit}
\index{TRIT!générateurs concrets}
\index{générateur elliptique}
\index{générateur parabolique}
\index{générateur hyperbolique}

Le \cref{chap:trit-algebras} fixe les trois types quadratiques abstraits
\(J^2=-I\), \(J^2=0\) et \(J^2=I\), tandis que le
\cref{thm:clasificacion-cuadratica-nonadica} fournit leurs modèles finis
sur \(\mathbb F_3\). Ces objets ne doivent pas être confondus avec les
réalisations réelles suivantes. On fixe ici trois espaces réels et trois
opérateurs concrets, chacun obtenu à partir de son antécédent causal déclaré.

Soit
\[
V_{\mathrm e}=A_2\otimes_{\mathbb Z}\mathbb R,
\qquad
V_{\mathrm p}=\mathbb R^2,
\qquad
V_{\mathrm h}=\mathbb R^2.
\]
Sur \(V_{\mathrm e}\) agit le Coxeter \(C\) de la
\cref{prop:coxeter-a2-primario}. Sur \(V_{\mathrm p}\), on fixe
\[
N=
\begin{pmatrix}
0&1\\
0&0
\end{pmatrix}.
\]
Enfin, le \cref{thm:descenso-necesario-fav} construit
\[
F_{\rm av}=
\begin{pmatrix}
0&1\\
1&1
\end{pmatrix},
\qquad
A:=F_{\rm av}^{\,2}=
\begin{pmatrix}
1&1\\
1&2
\end{pmatrix}
\]
sur \(V_{\mathrm h}\).

\begin{theorem}[Trois réalisations concrètes des régimes TRIT]
\label{thm:tres-generadores-concretos-trit}
Définissons
\[
J_{\mathrm e}:=\frac{2C+I}{\sqrt3},
\qquad
J_{\mathrm p}:=N,
\qquad
J_{\mathrm h}:=\frac{2A-3I}{\sqrt5}.
\]
Alors
\[
J_{\mathrm e}^{\,2}=-I_{V_{\mathrm e}},
\qquad
J_{\mathrm p}^{\,2}=0,
\qquad
J_{\mathrm h}^{\,2}=I_{V_{\mathrm h}}.
\]
Par conséquent, pour \(t\in\mathbb R\),
\[
\exp(tJ_{\mathrm e})
=\cos t\,I+\sin t\,J_{\mathrm e},
\qquad
\exp(tJ_{\mathrm p})
=I+tJ_{\mathrm p},
\]
\[
\exp(tJ_{\mathrm h})
=\cosh t\,I+\sinh t\,J_{\mathrm h}.
\]

De plus, dans la branche elliptique,
\[
C=-\frac12I+\frac{\sqrt3}{2}J_{\mathrm e},
\]
et, pour \(n\in\mathbb Z\),
\[
C^n=
\cos\!\left(\frac{2\pi n}{3}\right)I
+\sin\!\left(\frac{2\pi n}{3}\right)J_{\mathrm e}.
\]
Dans la branche hyperbolique, la forme
\[
G=
\begin{pmatrix}
2&1\\
1&-2
\end{pmatrix}
\]
a pour signature \((1,1)\) et satisfait
\[
A^{\mathsf T}GA=G,\qquad
\det A=1,\qquad
\operatorname{tr}A=3.
\]
Si
\[
\eta:=\operatorname{arcosh}\frac32=2\log\varphi,
\]
alors
\[
A=\cosh(\eta)I+\sinh(\eta)J_{\mathrm h},
\qquad
A^n=\cosh(n\eta)I+\sinh(n\eta)J_{\mathrm h}.
\]
\end{theorem}

\begin{proof}
Le polynôme caractéristique du Coxeter est \(X^2+X+1\) ; par
Cayley--Hamilton,
\[
C^2+C+I=0.
\]
D'où,
\[
(2C+I)^2=4C^2+4C+I=-3I,
\]
et donc \(J_{\mathrm e}^2=-I\). L'identité pour \(C^n\) découle de
\(C=-\frac12I+\frac{\sqrt3}{2}J_{\mathrm e}\).

La multiplication directe donne \(N^2=0\). La série exponentielle se
tronque donc exactement :
\[
\exp(tN)=I+tN.
\]

Pour la troisième branche,
\[
A^{\mathsf T}GA=G,\qquad
\det G=-5<0,
\]
de sorte que \(G\) est non dégénérée et possède la signature \((1,1)\). De même,
\[
(2A-3I)^2=5I,
\]
d'où \(J_{\mathrm h}^2=I\). Comme
\[
\cosh(2\log\varphi)=\frac32,
\qquad
\sinh(2\log\varphi)=\frac{\sqrt5}{2},
\]
on obtient
\[
A=\frac32I+\frac{\sqrt5}{2}J_{\mathrm h}
=\cosh(\eta)I+\sinh(\eta)J_{\mathrm h}.
\]
La formule pour \(A^n\) résulte de la loi exponentielle. Les trois expressions
générales s'obtiennent, dans chaque cas, en séparant dans la série exponentielle les
puissances paires et impaires.
\end{proof}

\begin{corollary}[Équation d'Euler étendue par le paramètre tritique]
\label{cor:euler-extendida-trit-fibonacci}
Pour \(\tau\in\{+1,0,-1\}\), soit \(J_\tau\) l'un des générateurs
précédents, normalisé par
\[
 J_\tau^2=-\tau I.
\]
Définissons les fonctions de la famille quadratique
\[
 C_\tau(t)=
 \begin{cases}
  \cos t,&\tau=+1,\\
  1,&\tau=0,\\
  \cosh t,&\tau=-1,
 \end{cases}
 \qquad
 S_\tau(t)=
 \begin{cases}
  \sin t,&\tau=+1,\\
  t,&\tau=0,\\
  \sinh t,&\tau=-1.
 \end{cases}
\]
Alors les trois branches satisfont une unique identité exponentielle :
\begin{equation}
 \boxed{\exp(tJ_\tau)=C_\tau(t)I+S_\tau(t)J_\tau.}
 \label{eq:euler-extendida-trit-fibonacci}
\end{equation}
En particulier,
\[
 \exp(\pi J_{+1})+I=0,
 \qquad
 \exp(J_0)=I+J_0,
 \qquad
 \exp\!\bigl(2\log\varphi\,J_{-1}\bigr)=F_{\rm av}^{\,2}.
\]
La première égalité est la fermeture matricielle d'Euler ; la deuxième réalise le
transport parabolique de seuil ; la troisième incorpore l'auto-échelle de
Fibonacci au moyen des valeurs propres \(\varphi\) et \(-\varphi^{-1}\).
\end{corollary}

\begin{proof}
L'identité \eqref{eq:euler-extendida-trit-fibonacci} s'obtient en séparant
les puissances paires et impaires de la série exponentielle et en utilisant
\(J_\tau^2=-\tau I\). Les trois spécialisations découlent de
\(J_{+1}^2=-I\), de \(J_0^2=0\) et de l'égalité déjà démontrée
\(F_{\rm av}^{\,2}=\exp(2\log\varphi\,J_{-1})\).
\end{proof}

\paragraph{\(3\) réalisations exactes des régimes quadratiques du TRIT}
Les branches elliptique, parabolique et hyperbolique réalisent, respectivement, le
Coxeter interne de \(A_2\), le générateur nilpotent et l'incidence
\(F_{\rm av}\) induite par le calendrier TPK. Les identités matricielles
sont exactes dans les espaces et calibres déclarés.

\paragraph{Conditions de réalisation des \(3\) régimes quadratiques du TRIT et obstructions de non-dégénérescence algébrique}
Le type TRIT isolé organise les trois possibilités quadratiques, mais ne
sélectionne pas à lui seul les opérateurs \(C\), \(N\) et \(F_{\rm av}\) : la
sélection utilise, respectivement, le Coxeter d'APP, le modèle nilpotent
de transport et la descente des modes de TPK. L'apparition de \(\pi\) dans la
formule spectrale elliptique appartient à l'écriture analytique du cycle et
ne remplace pas la génération régionale de \(\pi_{\rm HMT}\). L'identité
\(A^{\mathsf T}GA=G\) construit une transformation de Lorentz exacte en
dimension \(1+1\) ; elle n'identifie pas à elle seule cet écran à
l'espace-temps physique \(3+1\). Un échec de
\[
(2C+I)^2=-3I,\qquad N^2=0,\qquad
(2A-3I)^2=5I,\qquad A^{\mathsf T}GA=G
\]
réfuterait la réalisation correspondante.
